%\documentclass[12pt]{article}
%\usepackage{ctex} % 中文支持
%\usepackage{titling} % 标题格式定制
%\usepackage{lipsum} % 示例文本
%\usepackage{titlesec} % 标题格式定制
%\usepackage{graphicx}
%
%% 设置标题格式
%\titleformat{\section}{\centering\Large\bfseries}{\thesection}{1em}{}
%\titleformat{\subsection}{\centering\large\bfseries}{\thesubsection}{1em}{}
%\titleformat{\subsubsection}{\centering\normalsize\bfseries}{\thesubsubsection}{1em}{}
%
%% 调整标题位置
%\setlength{\droptitle}{-4cm}

\documentclass[UTF8]{ctexart}
\usepackage{lmodern}
\usepackage{amsmath}
\usepackage{amssymb}
\usepackage{graphicx}
\usepackage{geometry}
\usepackage{color}  %单元格纵向合并

\geometry{left=2.18cm,right=2.18cm,top=1.54cm,bottom=2.0cm}
\pagestyle{empty}
\title{\bf 2026春计算方法--实验报告 \#3}
\author{姓名：\underline{~~~~~~~~~~~~~~~~~~~~~~~}  学号：\underline{~~~~~~~~~~~~~~~~~~~~~~~} }
\date{\today}

\begin{document}

%\title{\huge{\textbf{2023秋计算方法-实验报告\#1}}}
%\author{姓名：\underline{   }\hspace{1cm}学号：PB2151xxxx}
\maketitle

运行环境：[自己给出。。。]
%win11，vscode，py3

\section*{实验内容与要求}

分别编写~Newton 迭代法 (通常也称 Newton-Raphson 迭代法)
$$\color{blue}x_{n+1}=x_n-\frac{f(x_n)}{f'(x_n)}$$
和 %Damped-Newton (DN) 迭代
霍氏迭代法
%\begin{center}
%\end{center}
   % $x_{k+1}=x_K-\tau\frac{f(x_k)}{f'(x_k)}$，(其中阻尼参数$\tau$: 0 <$\tau$< 1)
 $$\color{blue}x_{n+1}= x_n - \frac{2f(x_n)f'(x_n)}{2f'(x_n)f'(x_n) - f(x_n)f''(x_n)}, \, n=0,1,2,\cdots$$
的通用程序, 并利用它们对如下非线性方程
$$\color{blue}f(x)\triangleq{arctan(x)}+2sin(x) - 0.1958 =0$$
求根. %取阻尼参数$\tau$=0.55,
计算中的停止条件为~$\color{blue}|f(x_n)|<10^{-8}$~或迭代步数~$\color{blue}n>2\times 10^4$~(此时,可视为迭代失败).
提示：编程前分别手算出$\color{blue}~f'(x), f''(x)$。

\begin{itemize}
    \item 列表给出两种迭代方法在初始点 $x_0$ 依次取值~-200, -90, -25, -19, -13, -5, 0.1958, 5, 13, 25, 50, 80, 150~时的迭代步数
    （若迭代步数超过2万步，可视为迭代失败）
    以及相应的数值解~$\color{blue}x_n$~({\color{blue}保留小数点后6位}; 迭代失败时无需给出)；
    \item 比较并分析两种方法的优劣，给出合理的算法分析并作实验小结。
\end{itemize}
\clearpage

\section{数值结果（列表或作图）}
函数~$f(x)$~图像参见图1; 计算结果参见表1.
\begin{table}[htb]
\centering
\caption{实验结果: 精度~$\color{blue}\varepsilon = 10^{-8}$, 最大迭代步数为2万}
\begin{tabular}{cccccc}
\hline
初始点 & Newton法步数 & Newton数值解 & 霍氏法步数 & 霍氏数值解 \\
\hline
-200 & ? & ?& ? & ?   \\
-90 & ? & ? &? & ?  \\
$\cdots$  & $\cdots$ & &  $\cdots$ & $\cdots$  \\
$\vdots$  & $\vdots$ & & $\vdots$ & $\vdots$ \\
  &   &  &  &   \\
\hline
\end{tabular}
\end{table}

{\bf 针对表格，给出必要的说明或解释.}

\section{算法分析}
%从表1中，不难观察到
\begin{itemize}
    \item Blah blah blah  
    \item Blah blah blah  
    \item Blah blah blah  
\end{itemize}

%Blah blah blah

%\begin{table}[H]
%\centering
%\caption{收敛范围}
%\begin{tabular}{|c|c|c|}
%\hline
%列1 & 列2 & 列3 \\
%\hline
%初始点 & Newton & Damped-Newton\\
%0 & [-0.7, 1.7] & [-1.2, 1.8]\\
%\hline
%\end{tabular}
%\end{table}
%


\section{实验小结}
在本次数值实验中，我们分别测试了两种非线性方程求根的计算方法，
发现了Blah blah blah，
同时体会了Blah blah blah
%算法的选择的重要性。最终得到以下结果：

~Newton~迭代法的优点：Blah blah blah,
缺点是Blah blah blah;

霍氏迭代法的优点：Blah blah blah, 缺点是Blah blah blah.

\begin{itemize}
    \item 
    \item 
\end{itemize}


\end{document}
